\begin{corollary}[7.4.7]
\label{IV.2.s7.4.7-fr}
Supongamos que $\boldsymbol P$ satisface las condiciones
$(\mathrm P_{\mathrm I})$, $(\mathrm P'_{\mathrm I})$,
$(\mathrm P_{\mathrm{II}})$, $(\mathrm P_{\mathrm{III}})$ y
$(\mathrm P_{\mathrm{IV}})$. Entonces, si $A$ es un
$\boldsymbol P$-anillo (7.4.5), el morfismo canónico
$\operatorname{Spec}(A[[T_1,\ldots,T_r]])\to\operatorname{Spec}(A)$ es un
$\boldsymbol P$-morfismo. En particular, si además $A$ es íntegro, $K$
su cuerpo de fracciones, y si $\mathfrak p$ es un ideal primo de
$B=A[[T_1,\ldots,T_r]]$ tal que $\mathfrak p\cap A=0$, entonces se
cumple la propiedad $\boldsymbol P(B_{\mathfrak p},K)$.
\end{corollary}

\begin{proof}
En efecto, el morfismo canónico
$\operatorname{Spec}(B)\to\operatorname{Spec}(A)$ se factoriza como
\[
\operatorname{Spec}(A[[T_1,\ldots,T_r]])
\xrightarrow{f}\operatorname{Spec}(A[T_1,\ldots,T_r])
\xrightarrow{g}\operatorname{Spec}(A).
\]
Está claro que el morfismo $g$ es regular (\textbf 0, 17.3.7); en virtud de
(7.4.5), $A[T_1,\ldots,T_r]$ es un $\boldsymbol P$-anillo, de modo que
por (7.4.6) $f$ es un $\boldsymbol P$-morfismo. Como $g$ es regular y,
por tanto, también un $\boldsymbol P$-morfismo (7.3.5, (i)), lo mismo
vale para $g\circ f$ en virtud de $(\mathrm P'_{\mathrm I})$.

Se observará que la conclusión sigue siendo válida si, en lugar de suponer
que $\boldsymbol P$ satisface $(\mathrm P'_{\mathrm I})$ y que todo
morfismo regular es un $\boldsymbol P$-morfismo, se supone únicamente
que un morfismo compuesto $g\circ f$ es un $\boldsymbol P$-morfismo
cuando $g$ es regular y $f$ un $\boldsymbol P$-morfismo (condición en
cierto modo simétrica de $(\mathrm P'_{\mathrm I})$).
\end{proof}
